\documentclass{article}
\usepackage{amsmath, amssymb, geometry, booktabs, xcolor}
\geometry{a4paper, margin=1in}

\title{Homework Chapter 3: Introduction to Determinants}
\author{Tashfeen Omran}
\date{September 2026}

\begin{document}

\maketitle
\tableofcontents
\newpage

\section{Part 1: Comprehensive Introduction, Context, and Prerequisites}

\subsection{Core Concepts}
The \textbf{determinant} is a scalar value that can be computed from the elements of a square matrix and encodes certain properties of the linear transformation described by the matrix. For an $n \times n$ matrix $A$, the determinant is denoted by $\det A$ or $|A|$. 

Geometrically, the absolute value of the determinant represents the scaling factor by which a linear transformation changes areas (in 2D) or volumes (in 3D). The sign of the determinant indicates whether the transformation preserves or reverses the orientation of the space. Algebraically, a matrix is invertible if and only if its determinant is non-zero.

To compute determinants of matrices larger than $2 \times 2$, we use \textbf{cofactor expansion}. The $(i, j)$-cofactor of a matrix $A$, denoted $C_{ij}$, is defined as $C_{ij} = (-1)^{i+j} \det A_{ij}$, where $A_{ij}$ is the submatrix formed by deleting the $i$-th row and $j$-th column of $A$.

\subsection{Intuition and Derivation}
Consider a $2 \times 2$ matrix $A = \begin{bmatrix} a & b \\ c & d \end{bmatrix}$. The columns of $A$ can be viewed as vectors $\mathbf{u} = \begin{bmatrix} a \\ c \end{bmatrix}$ and $\mathbf{v} = \begin{bmatrix} b \\ d \end{bmatrix}$. The area of the parallelogram spanned by these vectors is given by $|ad - bc|$. If we apply row operations to reduce $A$ to an upper triangular matrix without swapping rows, the diagonal entries represent the base and height of this parallelogram. The product of these diagonal entries is exactly $ad - bc$. 

For larger matrices, the recursive cofactor expansion formula is derived by systematically applying row operations to reduce the matrix to a triangular form, keeping track of how each operation scales the volume of the $n$-dimensional parallelepiped.

\subsection{Historical Context and Motivation}
The concept of the determinant actually predates the modern concept of a matrix. In the late 17th century, mathematicians like Gottfried Wilhelm Leibniz and Seki Takakazu independently developed determinants to solve systems of linear equations. They wanted a single algebraic formula that could ``determine'' whether a system of equations had a unique solution. This led to Cramer's Rule in the 18th century. It wasn't until the 19th century that Arthur Cayley and James Joseph Sylvester formalized matrix algebra, placing the determinant in its modern context as a property of square matrices.

\subsection{Key Formulas}
\begin{itemize}
    \item \textbf{$1 \times 1$ Matrix:} If $A = [a]$, then $\det A = a$.
    \item \textbf{$2 \times 2$ Matrix:} If $A = \begin{bmatrix} a & b \\ c & d \end{bmatrix}$, then $\det A = ad - bc$.
    \item \textbf{Cofactor:} $C_{ij} = (-1)^{i+j} \det A_{ij}$.
    \item \textbf{Cofactor Expansion (across row $i$):} $\det A = a_{i1}C_{i1} + a_{i2}C_{i2} + \dots + a_{in}C_{in}$.
    \item \textbf{Cofactor Expansion (down column $j$):} $\det A = a_{1j}C_{1j} + a_{2j}C_{2j} + \dots + a_{nj}C_{nj}$.
    \item \textbf{Triangular Matrix:} If $A$ is triangular, $\det A = a_{11}a_{22} \dots a_{nn}$ (the product of the main diagonal entries).
\end{itemize}

\subsection{Prerequisites}
To master this topic, you must be comfortable with:
\begin{enumerate}
    \item \textbf{Matrix Terminology:} Understanding rows, columns, square matrices, and diagonal entries.
    \item \textbf{Elementary Row Operations:} Row replacement, row swapping, and scalar multiplication.
    \item \textbf{Basic Arithmetic:} Careful bookkeeping of negative signs is the most critical mechanical skill required here.
\end{enumerate}

\newpage
\section{Part 2: Detailed Homework Solutions}
\textit{Note: A strategic selection of 15 problems has been chosen to represent every distinct problem archetype found in Section 3.1, ensuring full conceptual coverage without redundant repetition.}

\subsection{Type 1: $3 \times 3$ Determinant Computation (Cofactor Expansion)}

\textbf{Exercise 1:} Compute the determinant using a cofactor expansion across the first row.
\[ A = \begin{bmatrix} 3 & 0 & 4 \\ 2 & 3 & 2 \\ 0 & 5 & -1 \end{bmatrix} \]

\textbf{Solution:}
We expand across the first row ($i=1$). The formula is $\det A = a_{11}C_{11} + a_{12}C_{12} + a_{13}C_{13}$.
\begin{align*}
\det A &= 3 \cdot (-1)^{1+1} \det \begin{bmatrix} 3 & 2 \\ 5 & -1 \end{bmatrix} + 0 \cdot (-1)^{1+2} \det \begin{bmatrix} 2 & 2 \\ 0 & -1 \end{bmatrix} + 4 \cdot (-1)^{1+3} \det \begin{bmatrix} 2 & 3 \\ 0 & 5 \end{bmatrix} \\
&= 3(1)((3)(-1) - (2)(5)) - 0 + 4(1)((2)(5) - (3)(0)) \\
&= 3(-3 - 10) + 4(10 - 0) \\
&= 3(-13) + 4(10) \\
&= -39 + 40 \\
&= 1
\end{align*}
\textbf{Final Answer:} $\det A = 1$

\vspace{0.5cm}
\textbf{Exercise 3:} Compute the determinant using a cofactor expansion across the first row.
\[ A = \begin{bmatrix} 2 & -2 & 3 \\ 3 & 1 & 2 \\ 1 & 3 & -1 \end{bmatrix} \]

\textbf{Solution:}
Expanding across the first row:
\begin{align*}
\det A &= 2 \det \begin{bmatrix} 1 & 2 \\ 3 & -1 \end{bmatrix} - (-2) \det \begin{bmatrix} 3 & 2 \\ 1 & -1 \end{bmatrix} + 3 \det \begin{bmatrix} 3 & 1 \\ 1 & 3 \end{bmatrix} \\
&= 2((1)(-1) - (2)(3)) + 2((3)(-1) - (2)(1)) + 3((3)(3) - (1)(1)) \\
&= 2(-1 - 6) + 2(-3 - 2) + 3(9 - 1) \\
&= 2(-7) + 2(-5) + 3(8) \\
&= -14 - 10 + 24 \\
&= 0
\end{align*}
\textbf{Final Answer:} $\det A = 0$

\subsection{Type 2: Exploiting Zeros in Larger Matrices}

\textbf{Exercise 11:} Compute the determinant by cofactor expansions. At each step, choose a row or column that involves the least amount of computation.
\[ A = \begin{bmatrix} 2 & -3 & 4 & 5 \\ 0 & 5 & 3 & -1 \\ 0 & 0 & -2 & 7 \\ 0 & 0 & 0 & 4 \end{bmatrix} \]

\textbf{Solution:}
Notice that matrix $A$ is an upper triangular matrix (all entries below the main diagonal are zero). According to Theorem 2, the determinant of a triangular matrix is simply the product of the entries on its main diagonal.
\begin{align*}
\det A &= a_{11} \cdot a_{22} \cdot a_{33} \cdot a_{44} \\
&= (2) \cdot (5) \cdot (-2) \cdot (4) \\
&= 10 \cdot (-8) \\
&= -80
\end{align*}
\textbf{Final Answer:} $\det A = -80$

\vspace{0.5cm}
\textbf{Exercise 13:} Compute the determinant.
\[ A = \begin{bmatrix} 4 & 0 & -7 & 3 & -5 \\ 0 & 0 & 2 & 0 & 0 \\ 7 & 3 & -6 & 4 & -8 \\ 5 & 0 & 5 & 2 & -3 \\ 0 & 0 & 9 & -1 & 2 \end{bmatrix} \]

\textbf{Solution:}
We seek the row or column with the most zeros. Column 2 has four zeros and only one non-zero entry ($a_{32} = 3$). We expand down Column 2. The sign for $a_{32}$ is $(-1)^{3+2} = -1$.
\[ \det A = -3 \det \begin{bmatrix} 4 & -7 & 3 & -5 \\ 0 & 2 & 0 & 0 \\ 5 & 5 & 2 & -3 \\ 0 & 9 & -1 & 2 \end{bmatrix} \]
Now, look at the new $4 \times 4$ matrix. Row 2 has three zeros and one non-zero entry ($2$ in the second column). We expand across Row 2. The sign for this entry (which is in the $(2,2)$ position of the submatrix) is $(-1)^{2+2} = +1$.
\[ \det A = -3 \left( 2 \det \begin{bmatrix} 4 & 3 & -5 \\ 5 & 2 & -3 \\ 0 & -1 & 2 \end{bmatrix} \right) = -6 \det \begin{bmatrix} 4 & 3 & -5 \\ 5 & 2 & -3 \\ 0 & -1 & 2 \end{bmatrix} \]
For the $3 \times 3$ matrix, we expand down Column 1 (which has one zero):
\begin{align*}
\det \begin{bmatrix} 4 & 3 & -5 \\ 5 & 2 & -3 \\ 0 & -1 & 2 \end{bmatrix} &= 4 \det \begin{bmatrix} 2 & -3 \\ -1 & 2 \end{bmatrix} - 5 \det \begin{bmatrix} 3 & -5 \\ -1 & 2 \end{bmatrix} + 0 \\
&= 4((2)(2) - (-3)(-1)) - 5((3)(2) - (-5)(-1)) \\
&= 4(4 - 3) - 5(6 - 5) \\
&= 4(1) - 5(1) \\
&= -1
\end{align*}
Finally, multiply by the constant we factored out:
\[ \det A = -6 \cdot (-1) = 6 \]
\textbf{Final Answer:} $\det A = 6$

\subsection{Type 3: The $3 \times 3$ Diagonal Trick}

\textbf{Exercise 15:} Compute the determinant using the downward/upward diagonal products trick.
\[ A = \begin{bmatrix} 1 & 0 & 4 \\ 2 & 3 & 2 \\ 0 & 5 & -2 \end{bmatrix} \]

\textbf{Solution:}
We copy the first two columns to the right of the matrix:
\[ \begin{matrix} 1 & 0 & 4 & | & 1 & 0 \\ 2 & 3 & 2 & | & 2 & 3 \\ 0 & 5 & -2 & | & 0 & 5 \end{matrix} \]
Sum of downward diagonals:
\[ D = (1)(3)(-2) + (0)(2)(0) + (4)(2)(5) = -6 + 0 + 40 = 34 \]
Sum of upward diagonals:
\[ U = (0)(3)(4) + (5)(2)(1) + (-2)(2)(0) = 0 + 10 + 0 = 10 \]
Determinant is $D - U$:
\[ \det A = 34 - 10 = 24 \]
\textbf{Final Answer:} $\det A = 24$

\vspace{0.5cm}
\textbf{Exercise 17:} Compute the determinant using the diagonal trick.
\[ A = \begin{bmatrix} 2 & -3 & 3 \\ 3 & 2 & 2 \\ 1 & 3 & -1 \end{bmatrix} \]

\textbf{Solution:}
Copy the first two columns:
\[ \begin{matrix} 2 & -3 & 3 & | & 2 & -3 \\ 3 & 2 & 2 & | & 3 & 2 \\ 1 & 3 & -1 & | & 1 & 3 \end{matrix} \]
Sum of downward diagonals:
\[ D = (2)(2)(-1) + (-3)(2)(1) + (3)(3)(3) = -4 - 6 + 27 = 17 \]
Sum of upward diagonals:
\[ U = (1)(2)(3) + (3)(2)(2) + (-1)(3)(-3) = 6 + 12 + 9 = 27 \]
Determinant is $D - U$:
\[ \det A = 17 - 27 = -10 \]
\textbf{Final Answer:} $\det A = -10$

\subsection{Type 4: Effect of Elementary Row Operations}

\textbf{Exercise 19:} Explore the effect of an elementary row operation on the determinant.
\[ \begin{bmatrix} a & b \\ c & d \end{bmatrix}, \quad \begin{bmatrix} a & b \\ c+ka & d+kb \end{bmatrix} \]

\textbf{Solution:}
Let $A = \begin{bmatrix} a & b \\ c & d \end{bmatrix}$. Its determinant is $\det A = ad - bc$.
Let $B = \begin{bmatrix} a & b \\ c+ka & d+kb \end{bmatrix}$. This matrix is obtained by adding $k$ times the first row to the second row (a row replacement operation).
We compute $\det B$:
\begin{align*}
\det B &= a(d+kb) - b(c+ka) \\
&= ad + akb - bc - abk \\
&= ad - bc
\end{align*}
\textbf{Final Answer:} $\det B = \det A$. A row replacement operation does not change the determinant.

\vspace{0.5cm}
\textbf{Exercise 21:} Explore the effect of an elementary row operation.
\[ \begin{bmatrix} 6 & 5 \\ 3 & 4 \end{bmatrix}, \quad \begin{bmatrix} 6 & 5 \\ 3+6k & 4+5k \end{bmatrix} \]

\textbf{Solution:}
Let $A = \begin{bmatrix} 6 & 5 \\ 3 & 4 \end{bmatrix}$. $\det A = (6)(4) - (5)(3) = 24 - 15 = 9$.
Let $B = \begin{bmatrix} 6 & 5 \\ 3+6k & 4+5k \end{bmatrix}$. This is obtained by adding $k$ times row 1 to row 2.
\begin{align*}
\det B &= 6(4+5k) - 5(3+6k) \\
&= 24 + 30k - 15 - 30k \\
&= 9
\end{align*}
\textbf{Final Answer:} $\det B = 9$. The determinant is unchanged by the row replacement operation.

\subsection{Type 5: Determinants of Elementary Matrices}

\textbf{Exercise 25:} Compute the determinant of the elementary matrix.
\[ E = \begin{bmatrix} 1 & 0 & 0 \\ 0 & 1 & 0 \\ 0 & k & 1 \end{bmatrix} \]

\textbf{Solution:}
This matrix is a lower triangular matrix. Its determinant is the product of its diagonal entries.
\[ \det E = 1 \cdot 1 \cdot 1 = 1 \]
\textbf{Final Answer:} $\det E = 1$

\vspace{0.5cm}
\textbf{Exercise 29:} Compute the determinant of the elementary matrix.
\[ E = \begin{bmatrix} 1 & 0 & 0 \\ 0 & k & 0 \\ 0 & 0 & 1 \end{bmatrix} \]

\textbf{Solution:}
This matrix is a diagonal matrix (which is a special case of a triangular matrix). Its determinant is the product of its diagonal entries.
\[ \det E = 1 \cdot k \cdot 1 = k \]
\textbf{Final Answer:} $\det E = k$

\subsection{Type 6: Theoretical Properties (True/False)}

\textbf{Exercise 39:} True or False: An $n \times n$ determinant is defined by determinants of $(n-1) \times (n-1)$ submatrices.

\textbf{Solution:}
\textbf{True.} The definition of the determinant for an $n \times n$ matrix (for $n \ge 2$) is given recursively as a sum of terms involving cofactors, where each cofactor $C_{ij}$ is defined in terms of the determinant of an $(n-1) \times (n-1)$ submatrix $A_{ij}$.

\vspace{0.5cm}
\textbf{Exercise 40:} True or False: The $(i, j)$-cofactor of a matrix $A$ is the matrix $A_{ij}$ obtained by deleting from $A$ its $i$-th row and $j$-th column.

\textbf{Solution:}
\textbf{False.} The matrix $A_{ij}$ is called a submatrix. The $(i, j)$-cofactor is a \textit{scalar number} defined as $C_{ij} = (-1)^{i+j} \det A_{ij}$. It is not a matrix itself.

\subsection{Type 7: Area and Volume Applications}

\textbf{Exercise 43:} Let $\mathbf{u} = \begin{bmatrix} 3 \\ 0 \end{bmatrix}$ and $\mathbf{v} = \begin{bmatrix} 1 \\ 2 \end{bmatrix}$. Compute the area of the parallelogram determined by $\mathbf{u}, \mathbf{v}, \mathbf{u}+\mathbf{v}$, and $\mathbf{0}$, and compute the determinant of $[\mathbf{u} \quad \mathbf{v}]$. How do they compare? Replace the first entry of $\mathbf{v}$ by an arbitrary number $x$, and repeat the problem.

\textbf{Solution:}
Let $A = [\mathbf{u} \quad \mathbf{v}] = \begin{bmatrix} 3 & 1 \\ 0 & 2 \end{bmatrix}$.
The determinant is $\det A = (3)(2) - (1)(0) = 6$.
The area of the parallelogram is the absolute value of the determinant: $|\det A| = |6| = 6$. They are exactly the same.

Now, replace $\mathbf{v}$ with $\mathbf{v}' = \begin{bmatrix} x \\ 2 \end{bmatrix}$.
Let $A' = [\mathbf{u} \quad \mathbf{v}'] = \begin{bmatrix} 3 & x \\ 0 & 2 \end{bmatrix}$.
The determinant is $\det A' = (3)(2) - (x)(0) = 6$.
The area is $|\det A'| = 6$.
\textbf{Final Answer:} The area and the determinant are equal (since the determinant is positive). Replacing the first entry of $\mathbf{v}$ with $x$ corresponds to a horizontal shear transformation, which shifts the top of the parallelogram parallel to the base but does not change its height or base length, thus preserving the area.

\subsection{Type 8: Bounding and Reasonable Answers}

\textbf{Exercise 45:} Let $A$ be a $2 \times 2$ matrix all of whose entries are numbers that are greater than or equal to $-10$ and less than or equal to $10$. Decide if each of the following is a reasonable answer for $\det A$.
a. $0$ \quad b. $202$ \quad c. $-110$ \quad d. $555$

\textbf{Solution:}
Let $A = \begin{bmatrix} a & b \\ c & d \end{bmatrix}$. We know $-10 \le a,b,c,d \le 10$.
The determinant is $\det A = ad - bc$.
To maximize $\det A$, we need to maximize $ad$ and minimize $bc$.
Max $ad = (10)(10) = 100$ or $(-10)(-10) = 100$.
Min $bc = (10)(-10) = -100$ or $(-10)(10) = -100$.
So, the maximum possible determinant is $100 - (-100) = 200$.
To minimize $\det A$, we minimize $ad$ and maximize $bc$.
Min $ad = -100$. Max $bc = 100$.
So, the minimum possible determinant is $-100 - 100 = -200$.
Therefore, the determinant must lie in the interval $[-200, 200]$.
\textbf{Final Answer:} 
a. $0$ is reasonable (in range).
b. $202$ is unreasonable (too large).
c. $-110$ is reasonable (in range).
d. $555$ is unreasonable (too large).

\subsection{Type 9: Design Matrix Application}

\textbf{Exercise 51:} Which of the following matrices corresponds to the best design for four weighings of four objects? The best design maximizes the determinant of $D^T D$.
a. $D = \begin{bmatrix} 1 & 1 & 1 & 1 \\ 1 & -1 & 1 & 1 \\ -1 & 1 & 1 & -1 \\ 1 & 1 & -1 & -1 \end{bmatrix}$
b. $D = \begin{bmatrix} -1 & -1 & -1 & -1 \\ 1 & -1 & 1 & -1 \\ 1 & 1 & -1 & -1 \\ 1 & 1 & 1 & -1 \end{bmatrix}$
c. $D = \begin{bmatrix} 1 & 1 & 1 & 1 \\ 1 & -1 & -1 & 1 \\ 1 & 1 & -1 & -1 \\ 1 & -1 & 1 & -1 \end{bmatrix}$

\textbf{Solution:}
We know from properties of determinants that $\det(D^T D) = \det(D^T)\det(D) = (\det D)^2$. Thus, we just need to find which matrix has the largest absolute determinant.

For matrix (a):
\[ \det \begin{bmatrix} 1 & 1 & 1 & 1 \\ 1 & -1 & 1 & 1 \\ -1 & 1 & 1 & -1 \\ 1 & 1 & -1 & -1 \end{bmatrix} \]
Apply $R_2 \leftarrow R_2 - R_1$, $R_3 \leftarrow R_3 + R_1$, $R_4 \leftarrow R_4 - R_1$:
\[ \det \begin{bmatrix} 1 & 1 & 1 & 1 \\ 0 & -2 & 0 & 0 \\ 0 & 2 & 2 & 0 \\ 0 & 0 & -2 & -2 \end{bmatrix} = 1 \cdot \det \begin{bmatrix} -2 & 0 & 0 \\ 2 & 2 & 0 \\ 0 & -2 & -2 \end{bmatrix} \]
This is a lower triangular matrix. The determinant is $(-2)(2)(-2) = 8$.
$(\det D)^2 = 64$.

For matrix (b):
\[ \det \begin{bmatrix} -1 & -1 & -1 & -1 \\ 1 & -1 & 1 & -1 \\ 1 & 1 & -1 & -1 \\ 1 & 1 & 1 & -1 \end{bmatrix} \]
Apply $R_1 \leftarrow R_1 + R_4$, $R_2 \leftarrow R_2 - R_4$, $R_3 \leftarrow R_3 - R_4$:
\[ \det \begin{bmatrix} 0 & 0 & 0 & -2 \\ 0 & -2 & 0 & 0 \\ 0 & 0 & -2 & 0 \\ 1 & 1 & 1 & -1 \end{bmatrix} \]
Expand across Row 1:
\[ -(-2) \det \begin{bmatrix} 0 & -2 & 0 \\ 0 & 0 & -2 \\ 1 & 1 & 1 \end{bmatrix} = 2 \left( 1 \cdot \det \begin{bmatrix} -2 & 0 \\ 0 & -2 \end{bmatrix} \right) = 2(4) = 8 \]
$(\det D)^2 = 64$.

For matrix (c):
\[ \det \begin{bmatrix} 1 & 1 & 1 & 1 \\ 1 & -1 & -1 & 1 \\ 1 & 1 & -1 & -1 \\ 1 & -1 & 1 & -1 \end{bmatrix} \]
Apply $R_2 \leftarrow R_2 - R_1$, $R_3 \leftarrow R_3 - R_1$, $R_4 \leftarrow R_4 - R_1$:
\[ \det \begin{bmatrix} 1 & 1 & 1 & 1 \\ 0 & -2 & -2 & 0 \\ 0 & 0 & -2 & -2 \\ 0 & -2 & 0 & -2 \end{bmatrix} = 1 \cdot \det \begin{bmatrix} -2 & -2 & 0 \\ 0 & -2 & -2 \\ -2 & 0 & -2 \end{bmatrix} \]
Expand across Column 1:
\[ -2 \det \begin{bmatrix} -2 & -2 \\ 0 & -2 \end{bmatrix} - 0 + (-2) \det \begin{bmatrix} -2 & 0 \\ -2 & -2 \end{bmatrix} = -2(4) - 2(4) = -8 - 8 = -16 \]
$(\det D)^2 = (-16)^2 = 256$.

\textbf{Final Answer:} Matrix (c) provides the best design because it maximizes $\det(D^T D) = 256$. In this design, for the first weighing, all objects are placed in the left pan. For the second weighing, objects $s_1, s_4$ are in the left pan and $s_2, s_3$ are in the right pan, and so on.

\newpage
\section{Part 3: In-Depth Analysis of Problems and Techniques}

\subsection{A) Problem Types and General Approach}
\begin{itemize}
    \item \textbf{Direct Computation (Cofactor Expansion):} Problems like Ex 1 and 3 require brute-force application of the recursive formula. The strategy is to pick a row or column, apply the checkerboard sign pattern, and compute the smaller determinants.
    \item \textbf{Exploiting Zeros:} Problems like Ex 11 and 13 test your ability to work smart, not hard. The strategy is to scan the matrix for the row or column with the maximum number of zeros. This collapses a massive sum into a single term.
    \item \textbf{Triangular Matrices:} A sub-type of exploiting zeros. If a matrix is upper or lower triangular, the determinant is instantly found by multiplying the main diagonal.
    \item \textbf{Diagonal Trick ($3 \times 3$ only):} Problems like Ex 15 and 17 use Sarrus' Rule. The strategy is to copy the first two columns, sum the products of the three downward diagonals, and subtract the sum of the products of the three upward diagonals.
    \item \textbf{Row Operation Analysis:} Problems like Ex 19 and 21 bridge computation and theory. The strategy is to compute the determinant before and after a row operation to empirically verify Theorem 3 (e.g., row replacement preserves the determinant).
\end{itemize}

\subsection{B) Key Algebraic and Calculus Manipulations}
\begin{itemize}
    \item \textbf{The Checkerboard Sign Pattern:} The most crucial algebraic manipulation is applying $(-1)^{i+j}$. For example, in Ex 13, expanding down column 2 meant the entry $a_{32}$ carried a negative sign because $3+2=5$ (odd). Forgetting this flips the sign of the entire answer.
    \item \textbf{Factoring out Constants:} When expanding, you pull the matrix entry out as a scalar multiplier for the sub-determinant. In Ex 13, we pulled out $-3$, then $2$, creating a running product of $-6$ outside the remaining $3 \times 3$ matrix.
    \item \textbf{Strategic Row Operations for Simplification:} In Ex 51, computing the determinant of a dense $4 \times 4$ matrix directly would require $4! = 24$ multiplications. Instead, we used row replacement operations (which do not change the determinant) to create zeros, effectively turning a dense matrix into a sparse one before expanding.
\end{itemize}

\section{Part 4: "Cheatsheet" and Tips for Success}

\subsection{Summary of Important Formulas}
\begin{itemize}
    \item \textbf{$2 \times 2$ Determinant:} $\det \begin{bmatrix} a & b \\ c & d \end{bmatrix} = ad - bc$
    \item \textbf{Cofactor:} $C_{ij} = (-1)^{i+j} \det A_{ij}$
    \item \textbf{Triangular Matrix Determinant:} $\det A = \prod_{i=1}^n a_{ii}$
\end{itemize}

\subsection{Tricks and Shortcuts}
\begin{itemize}
    \item \textbf{The Zero Hunter:} Always spend 5 seconds scanning a matrix before calculating. If a row or column is entirely zeros, the determinant is $0$. If a row has mostly zeros, expand along it.
    \item \textbf{The Sarrus Shortcut:} For $3 \times 3$ matrices, the diagonal trick is usually faster and less prone to sign errors than cofactor expansion.
    \item \textbf{Row Replacement Cheat:} If a matrix is dense (no zeros), use the operation $R_i \leftarrow R_i + kR_j$ to create zeros. This operation is "free" because it does not change the determinant's value.
\end{itemize}

\subsection{Common Pitfalls and Traps}
\begin{itemize}
    \item \textbf{The $4 \times 4$ Diagonal Trap:} Sarrus' Rule (the diagonal trick) \textbf{ONLY} works for $3 \times 3$ matrices. Applying it to a $4 \times 4$ matrix is a fatal error that yields complete nonsense.
    \item \textbf{Sign Errors in Cofactors:} Forgetting that the $(1,2)$ position carries a negative sign is the \#1 reason students get these problems wrong. Always draw a quick $+ - +$ checkerboard in the margins.
    \item \textbf{Row Scaling Trap:} While row \textit{replacement} doesn't change the determinant, multiplying a row by a scalar $k$ multiplies the determinant by $k$. Do not confuse the two operations!
\end{itemize}

\newpage
\section{Part 5: Conceptual Synthesis and The "Big Picture"}

\subsection{A) Thematic Connections}
The core theme of the determinant is \textbf{volume scaling and invertibility}. In algebra, we often want to know if an operation can be undone (inverted). For real numbers, any number $x$ has a multiplicative inverse $1/x$ as long as $x \neq 0$. The determinant generalizes this concept to higher dimensions: a matrix transformation can be undone as long as it doesn't squash the space into a lower dimension (i.e., as long as its volume scaling factor, the determinant, is not zero).

\subsection{B) Forward and Backward Links}
\begin{itemize}
    \item \textbf{Backward Links:} This topic builds directly on the elementary row operations learned in Chapter 1. The very operations used to find Reduced Row Echelon Form (RREF) are the ones that manipulate the determinant.
    \item \textbf{Forward Links:} The determinant is the absolute bedrock for finding \textbf{Eigenvalues} in Chapter 5. To find the eigenvalues of $A$, you must solve the characteristic equation $\det(A - \lambda I) = 0$. Furthermore, in Multivariable Calculus, the determinant of the Jacobian matrix is required to change variables in multiple integrals, ensuring that the "infinitesimal volume" is scaled correctly.
\end{itemize}

\section{Part 6: Abstract Real-World Modeling (AI, ML, \& Quant Finance)}

\subsection{A) Concrete Scenarios}
\begin{enumerate}
    \item \textbf{Quantitative Finance (Portfolio Optimization):} In Modern Portfolio Theory, the covariance matrix $\Sigma$ captures the variance and correlations of asset returns. The determinant of $\Sigma$, known as the \textit{generalized variance}, measures the overall multidimensional volatility of the portfolio. If $\det \Sigma$ is close to zero, it indicates severe multicollinearity (assets are perfectly correlated), meaning the portfolio is not truly diversified and the inverse matrix $\Sigma^{-1}$ (needed for optimal weight allocation) will be highly unstable.
    \item \textbf{Generative AI (Normalizing Flows):} In deep learning, Normalizing Flows are used to model complex probability distributions (like images of faces). They work by passing a simple distribution (like a standard Gaussian) through a series of invertible neural network layers. To compute the exact likelihood of the generated data, the model must calculate the determinant of the Jacobian matrix of these transformations.
\end{enumerate}

\subsection{B) Model Problem Setup: Normalizing Flows in Machine Learning}
Let $\mathbf{z} \in \mathbb{R}^n$ be a latent random variable drawn from a simple base distribution $p_Z(\mathbf{z}) = \mathcal{N}(\mathbf{0}, I)$. We wish to model a complex, high-dimensional data distribution $p_X(\mathbf{x})$ (e.g., natural images). 

We define a bijective, differentiable transformation $f_\theta : \mathbb{R}^n \to \mathbb{R}^n$ parameterized by a neural network with weights $\theta$, such that $\mathbf{x} = f_\theta(\mathbf{z})$. By the change of variables formula from probability theory, the exact log-likelihood of a data point $\mathbf{x}$ is given by:
\[ \log p_X(\mathbf{x}) = \log p_Z(f_\theta^{-1}(\mathbf{x})) + \log \left| \det J_{f_\theta^{-1}}(\mathbf{x}) \right| \]
where $J_{f_\theta^{-1}}$ is the $n \times n$ Jacobian matrix of the inverse transformation.

\textbf{The Computational Bottleneck:} Computing the determinant of an arbitrary $n \times n$ matrix takes $\mathcal{O}(n^3)$ time. For high-resolution images, $n$ can be $10^6$, making $\mathcal{O}(n^3)$ impossible. 

\textbf{The Solution (Using Theorem 2):} AI researchers design the neural network $f_\theta$ such that its Jacobian matrix is strictly \textit{lower triangular}. Because the determinant of a triangular matrix is simply the product of its diagonal entries (Theorem 2 from this chapter), the computation drops from $\mathcal{O}(n^3)$ to $\mathcal{O}(n)$. The loss function to be minimized via stochastic gradient descent becomes:
\[ \mathcal{L}(\theta) = - \mathbb{E}_{\mathbf{x} \sim \text{data}} \left[ \log p_Z(f_\theta^{-1}(\mathbf{x})) + \sum_{i=1}^n \log \left| \frac{\partial (f_\theta^{-1})_i}{\partial x_i} \right| \right] \]
This elegant application of triangular determinants makes exact likelihood training of deep generative models computationally tractable.

\newpage
\section{Part 7: Common Variations and Untested Concepts}

\subsection{Missing Concept 1: The Vandermonde Determinant}
A classic determinant problem not covered in this specific exercise set is the Vandermonde matrix, which appears frequently in polynomial interpolation.
A $3 \times 3$ Vandermonde matrix looks like this:
\[ V = \begin{bmatrix} 1 & x_1 & x_1^2 \\ 1 & x_2 & x_2^2 \\ 1 & x_3 & x_3^2 \end{bmatrix} \]
\textbf{Example:} Prove that $\det V = (x_2 - x_1)(x_3 - x_1)(x_3 - x_2)$.
\textit{Brief Solution:} Apply row operations $R_2 \leftarrow R_2 - R_1$ and $R_3 \leftarrow R_3 - R_1$:
\[ \det V = \det \begin{bmatrix} 1 & x_1 & x_1^2 \\ 0 & x_2 - x_1 & x_2^2 - x_1^2 \\ 0 & x_3 - x_1 & x_3^2 - x_1^2 \end{bmatrix} \]
Expand down column 1, factor out $(x_2 - x_1)$ from the new row 1 and $(x_3 - x_1)$ from the new row 2, and the remaining $2 \times 2$ determinant evaluates to $(x_3 - x_2)$.

\subsection{Missing Concept 2: Block Matrix Determinants}
Often, large matrices can be partitioned into smaller blocks. If a matrix $M$ is block upper-triangular, $M = \begin{bmatrix} A & B \\ 0 & D \end{bmatrix}$, where $A$ and $D$ are square matrices, then $\det M = (\det A)(\det D)$. This is a powerful shortcut for matrices with large blocks of zeros.

\section{Part 8: Advanced Diagnostic Testing: "Find the Flaw"}

\textit{Instructions: Act as the grader. For each problem, identify the specific error, explain why it is an error in one sentence, and provide the correct solution.}

\vspace{0.3cm}
\textbf{Problem 1:} Compute the determinant of $A = \begin{bmatrix} 1 & 0 & 0 & 2 \\ 0 & 3 & 0 & 0 \\ 0 & 0 & 4 & 0 \\ 5 & 0 & 0 & 6 \end{bmatrix}$.
\newline \textit{Flawed Solution:} Using the diagonal trick, the downward diagonal is $(1)(3)(4)(6) = 72$. The upward diagonal is $(5)(0)(0)(2) = 0$. Therefore, $\det A = 72 - 0 = 72$.

\textbf{Problem 2:} Compute the determinant of $A = \begin{bmatrix} 2 & 0 & 1 \\ 3 & 0 & 4 \\ 1 & 5 & 2 \end{bmatrix}$.
\newline \textit{Flawed Solution:} Expand down the second column. The only non-zero entry is $5$.
$\det A = 5 \cdot \det \begin{bmatrix} 2 & 1 \\ 3 & 4 \end{bmatrix} = 5(8 - 3) = 5(5) = 25$.

\textbf{Problem 3:} Let $A$ be a $3 \times 3$ matrix with $\det A = 5$. Find $\det(2A)$.
\newline \textit{Flawed Solution:} Since the determinant is a linear operator, we can pull the constant out. $\det(2A) = 2 \det A = 2(5) = 10$.

\textbf{Problem 4:} Compute the determinant of $A = \begin{bmatrix} 0 & 2 \\ 3 & 4 \end{bmatrix}$.
\newline \textit{Flawed Solution:} Swap row 1 and row 2 to get an upper triangular matrix $B = \begin{bmatrix} 3 & 4 \\ 0 & 2 \end{bmatrix}$. The determinant of a triangular matrix is the product of its diagonals, so $\det A = (3)(2) = 6$.

\textbf{Problem 5:} Let $A = \begin{bmatrix} 1 & 2 \\ 0 & 3 \end{bmatrix}$ and $B = \begin{bmatrix} 4 & 0 \\ 1 & 2 \end{bmatrix}$. Find $\det(A + B)$.
\newline \textit{Flawed Solution:} $\det A = 3$ and $\det B = 8$. Therefore, $\det(A + B) = \det A + \det B = 3 + 8 = 11$.

\vspace{0.5cm}
\hrule
\vspace{0.5cm}

\textbf{Grader's Answer Key:}

\textbf{Problem 1 Correction:}
\textit{Error:} The student applied Sarrus' Rule (the diagonal trick) to a $4 \times 4$ matrix, which is mathematically invalid.
\textit{Correct Solution:} Expand across row 2: $\det A = 3 \det \begin{bmatrix} 1 & 0 & 2 \\ 0 & 4 & 0 \\ 5 & 0 & 6 \end{bmatrix}$. Then expand across row 2 of the submatrix: $3(4 \det \begin{bmatrix} 1 & 2 \\ 5 & 6 \end{bmatrix}) = 12(6 - 10) = 12(-4) = -48$.

\textbf{Problem 2 Correction:}
\textit{Error:} The student forgot to apply the checkerboard sign pattern; the entry $5$ is in the $(3,2)$ position, so its sign should be $(-1)^{3+2} = -1$.
\textit{Correct Solution:} $\det A = -5 \cdot \det \begin{bmatrix} 2 & 1 \\ 3 & 4 \end{bmatrix} = -5(8 - 3) = -5(5) = -25$.

\textbf{Problem 3 Correction:}
\textit{Error:} The student assumed $\det(kA) = k \det A$, but scaling an $n \times n$ matrix scales every row, so the constant is pulled out $n$ times.
\textit{Correct Solution:} For an $n \times n$ matrix, $\det(kA) = k^n \det A$. Since $A$ is $3 \times 3$, $\det(2A) = 2^3 \det A = 8(5) = 40$.

\textbf{Problem 4 Correction:}
\textit{Error:} The student forgot that swapping two rows multiplies the determinant by $-1$.
\textit{Correct Solution:} Swapping the rows gives $B$, where $\det B = 6$. Since $B$ was obtained by one row swap, $\det A = -\det B = -6$. (Alternatively, directly compute $(0)(4) - (2)(3) = -6$).

\textbf{Problem 5 Correction:}
\textit{Error:} The student assumed the determinant is linear over matrix addition ($\det(A+B) = \det A + \det B$), which is generally false.
\textit{Correct Solution:} First, compute the matrix sum: $A + B = \begin{bmatrix} 5 & 2 \\ 1 & 5 \end{bmatrix}$. Then compute the determinant: $\det(A+B) = (5)(5) - (2)(1) = 25 - 2 = 23$.

\end{document}